
\documentclass{article}
\usepackage{colortbl}
\usepackage{makecell}
\usepackage{multirow}
\usepackage{supertabular}

\begin{document}

\newcounter{utterance}

\centering \large Interaction Transcript for game `cladder', experiment `full\_v1.5\_default', episode 7123 with Qwen3.5{-}4B.
\vspace{24pt}

{ \footnotesize  \setcounter{utterance}{1}
\setlength{\tabcolsep}{0pt}
\begin{supertabular}{c@{$\;$}|p{.15\linewidth}@{}p{.15\linewidth}p{.15\linewidth}p{.15\linewidth}p{.15\linewidth}p{.15\linewidth}}
    \# & \multicolumn{2}{c}{Player} && \multicolumn{2}{c}{Game Master} \\
    \hline

    \theutterance \stepcounter{utterance}  
    & & & \multicolumn{4}{p{0.6\linewidth}}{
        \cellcolor[rgb]{0.9,0.9,0.9}{
            \makecell[{{p{\linewidth}}}]{
                \texttt{\tiny{[P1$\langle$GM]}}
                \texttt{You are an expert in causal inference. The following question is not a typical commonsense query, but rather a meticulously designed question created by a professor specializing in causal inference, intended to assess the students' mastery of the course content.} \\
\texttt{Imagine a self{-}contained, hypothetical world with only the following conditions, and without any unmentioned factors or causal relationships: Jyka has a direct effect on yupt and kwox. Yupt has a direct effect on kwox. Method 1: We look at how jyka correlates with kwox case by case according to yupt. Method 2: We look directly at how jyka correlates with kwox in general. To understand how jyka affects kwox, is it more correct to use the Method 1 than Method 2?} \\
\texttt{Start your answer with 'Yes' or 'No', followed by additional reasoning or evidence to support your explanation.} \\
            }
        }
    }
     \\ \\

    \theutterance \stepcounter{utterance}  
    & \multicolumn{4}{p{0.6\linewidth}}{
        \cellcolor[rgb]{0.561,0.737,0.561}{
            \makecell[{{p{\linewidth}}}]{
                \texttt{\tiny{[P1$\rangle$GM]}}
                \texttt{Yes} \\
\\ 
\texttt{**Reasoning:**} \\
\\ 
\texttt{In this hypothetical world, we have the following causal structure:} \\
\texttt{1.  \$Jyka \textbackslash{}rightarrow Kwox\$ (Direct effect)} \\
\texttt{2.  \$Jyka \textbackslash{}rightarrow Yupt\$ (Direct effect)} \\
\texttt{3.  \$Yupt \textbackslash{}rightarrow Kwox\$ (Direct effect)} \\
\\ 
\texttt{This forms a classic **collider** structure where \$Jyka\$ and \$Yupt\$ both point to \$Kwox\$ (i.e., \$Jyka \textbackslash{}rightarrow Kwox \textbackslash{}leftarrow Yupt\$).} \\
\\ 
\texttt{*   **Analysis of Method 2 (General Correlation):** If we simply look at the correlation between \$Jyka\$ and \$Kwox\$ without conditioning on \$Yupt\$, we are observing the total effect of \$Jyka\$ on \$Kwox\$. However, because \$Kwox\$ is a collider, the path from \$Jyka\$ to \$Kwox\$ is independent of the path from \$Yupt\$ to \$Kwox\$ in the absence of conditioning. Therefore, the marginal correlation in Method 2 reflects only the direct causal effect of \$Jyka\$ on \$Kwox\$ (plus any noise). While this tells us the direct effect, it does not isolate the mechanism or account for the specific pathway through \$Yupt\$ if the goal is to understand the *total* impact including indirect routes, or if the question implies isolating the specific causal channel. More critically, in causal inference, "understanding how A affects B" often requires distinguishing between direct and indirect effects or ensuring we are not inadvertently inducing spurious associations if we were to condition on the collider.} \\
\\ 
\texttt{*   **Analysis of Method 1 (Conditioning on Yupt):** Method 1 involves looking at the relationship between \$Jyka\$ and \$Kwox\$ *case by case according to \$Yupt\$*. This is equivalent to conditioning on the variable \$Yupt\$ (stratification or regression controlling for \$Yupt\$).} \\
\texttt{    *   By conditioning on \$Yupt\$, we block the path \$Jyka \textbackslash{}rightarrow Yupt \textbackslash{}rightarrow Kwox\$.} \\
\texttt{    *   This isolates the **direct causal effect** of \$Jyka\$ on \$Kwox\$.} \\
\texttt{    *   In causal diagrams, conditioning on a non{-}collider (like \$Yupt\$) on a path blocks that path. Since \$Yupt\$ is a parent of the collider \$Kwox\$, conditioning on it does not open a backdoor path between \$Jyka\$ and \$Kwox\$ (there are no other paths here).} \\
\texttt{    *   Therefore, Method 1 provides a clean estimate of the specific causal effect of \$Jyka\$ on \$Kwox\$ that is not mediated by \$Yupt\$.} \\
\\ 
\texttt{**Conclusion:**} \\
\texttt{Method 1 is more correct for understanding the specific causal influence of \$Jyka\$} \\
            }
        }
    }
    & & \\ \\

    \theutterance \stepcounter{utterance}  
    & & & \multicolumn{2}{p{0.3\linewidth}}{
        \cellcolor[rgb]{0.9,0.9,0.9}{
            \makecell[{{p{\linewidth}}}]{
                \texttt{\tiny{[GM$|$GM]}}
                \texttt{Yes} \\
            }
        }
    }
    & & \\ \\

    \theutterance \stepcounter{utterance}  
    & & & \multicolumn{2}{p{0.3\linewidth}}{
        \cellcolor[rgb]{0.9,0.9,0.9}{
            \makecell[{{p{\linewidth}}}]{
                \texttt{\tiny{[GM$|$GM]}}
                \texttt{no} \\
            }
        }
    }
    & & \\ \\

    \theutterance \stepcounter{utterance}  
    & & & \multicolumn{2}{p{0.3\linewidth}}{
        \cellcolor[rgb]{0.9,0.9,0.9}{
            \makecell[{{p{\linewidth}}}]{
                \texttt{\tiny{[GM$|$GM]}}
                \texttt{game\_result = LOSE} \\
            }
        }
    }
    & & \\ \\

\end{supertabular}
}

\end{document}
